\section*{Hoofdstuk \ref{ch:g6:rocks-and-minerals} --- Gesteenten en mineralen}

\begin{solution}{exo:g6:rocks-and-minerals:1}
Graniet: \omterm{def:g6:rocks-and-minerals:rock-families}{stollingsgesteente}. Zandsteen: \omterm{def:g6:rocks-and-minerals:rock-families}{afzettingsgesteente}. Marmer:
metamorf. Basalt: \omterm{def:g6:rocks-and-minerals:rock-families}{stollingsgesteente}. Kalksteen: \omterm{def:g6:rocks-and-minerals:rock-families}{afzettingsgesteente}.
Leisteen: metamorf.
\end{solution}

\begin{solution}{exo:g6:rocks-and-minerals:2}
Kwarts is een \omterm{def:g6:rocks-and-minerals:mineral}{mineraal}; graniet is een \omterm{def:g3:stones-pebbles-sand:rock}{gesteente}. Een \omterm{def:g6:rocks-and-minerals:mineral}{mineraal} is één
natuurlijke stof met atomen die regelmatig geordend zijn; een \omterm{def:g3:stones-pebbles-sand:rock}{gesteente}
bestaat uit één of meer \omterm{def:g6:rocks-and-minerals:mineral}{mineralen} die aan elkaar vastzitten (graniet:
veldspaat, kwarts en mica).
\end{solution}

\begin{solution}{exo:g6:rocks-and-minerals:3}
Talk (1), calciet (3), kwarts (7), diamant (10). Diamant, het hardste,
krast alle andere.
\end{solution}

\begin{solution}{exo:g6:rocks-and-minerals:4}
``Gesmolten \omterm{def:g3:stones-pebbles-sand:rock}{gesteente} wordt een \omterm{def:g6:rocks-and-minerals:rock-families}{stollingsgesteente} door afkoelen. Een
kalksteen wordt een marmer door warmte en druk. Een graniet valt uiteen in
zandkorrels door \omterm{def:g5:water-shapes-the-land:weathering}{verwering}.''
\end{solution}

\begin{solution}{exo:g6:rocks-and-minerals:5}
(a) Basalt. (b) Marmer. (c) Zandsteen.
\end{solution}

\begin{solution}{exo:g6:rocks-and-minerals:6}
De \omterm{def:g6:rocks-and-minerals:hardness}{hardheid} is groter dan die van een vingernagel (ongeveer 2.5) en kleiner
dan die van een koperen munt (ongeveer 3.5). Het zou calciet (3) kunnen
zijn.
\end{solution}

\begin{solution}{exo:g6:rocks-and-minerals:7}
Kwarts: $21 \div 60 = 0.35$, dat is 35~\%. Veldspaat: $33 \div 60 = 0.55$,
dat is 55~\%. Mica: $6 \div 60 = 0.10$, dat is 10~\%. Controle:
$35 + 55 + 10 = 100$.
\end{solution}

\begin{solution}{exo:g6:rocks-and-minerals:8}
Graniet koelt heel langzaam af, diep onder de grond: de korrels hebben veel
tijd om te groeien en worden groot. Basalt komt uit \omterm{def:g4:volcanoes-earthquakes-first-look:eruption}{lava}, die snel afkoelt
aan het oppervlak: de korrels hebben geen tijd om te groeien en blijven te
klein om te zien.
\end{solution}

\begin{solution}{exo:g6:rocks-and-minerals:9}
(a) Kalksteen is een \omterm{def:g6:rocks-and-minerals:rock-families}{afzettingsgesteente}; marmer is metamorf. (b) Diep
onder de grond hebben warmte en druk de kalksteen omgevormd zonder hem te
laten smelten: het calciet groeide uit tot glimmende kristallen die in
elkaar passen. (c) Allebei bestaan ze uit hetzelfde \omterm{def:g6:rocks-and-minerals:mineral}{mineraal}, calciet, en
het is calciet dat bruist onder zuur.
\end{solution}

\begin{solution}{exo:g6:rocks-and-minerals:10}
(a) Ja: kwarts (7) is harder dan glas (ongeveer 5.5). (b) Nee: fluoriet (4)
is zachter dan glas. (c) Ja: een stalen spijker (ongeveer 6.5) is harder
dan veldspaat (6).
\end{solution}

\begin{solution}{exo:g6:rocks-and-minerals:11}
Bijvoorbeeld: (1) het graniet van de \omterm{def:g1:landscapes-around-us:landform}{berg} (\omterm{def:g6:rocks-and-minerals:rock-families}{stollingsgesteente}) valt uiteen
door \omterm{def:g5:water-shapes-the-land:weathering}{verwering}, en de kwartskorrels komen vrij; (2) een \omterm{def:g1:landscapes-around-us:water}{rivier} voert ze
weg: \omterm{def:g5:water-shapes-the-land:erosion}{erosie}; (3) de korrels worden in lagen op de bodem van de \omterm{def:g2:water-on-earth:ocean}{zee}
neergelegd: \omterm{def:g5:water-shapes-the-land:deposition}{afzetting}; (4) begraven onder andere lagen worden ze
samengeperst en gekit tot een zandsteen: verharding (het \omterm{def:g3:stones-pebbles-sand:rock}{gesteente} is nu
een \omterm{def:g6:rocks-and-minerals:rock-families}{afzettingsgesteente}); (5) diep onder de grond terechtgekomen wordt het
zandsteen door warmte en druk omgevormd tot een \omterm{def:g6:rocks-and-minerals:rock-families}{metamorf gesteente}.
\end{solution}

\begin{solution}{exo:g6:rocks-and-minerals:12}
(a) Niet gekrast door staal (ongeveer 6.5) en in staat glas te krassen:
harder dan 6.5. Gekrast door topaas (8): zachter dan 8. De \omterm{def:g6:rocks-and-minerals:hardness}{hardheid} is 7:
het is kwarts. (b) Kwarts is harder dan bijna alles wat ertegen wrijft
(andere korrels, \omterm{def:g3:stones-pebbles-sand:pebble}{kiezels}): het slijt nauwelijks af. Calciet (3) en mica
(2 tot 3) zijn veel zachter: door het wrijven worden ze vermalen tot
piepkleine stukjes die worden weggevoerd.
\end{solution}

\begin{solution}{pb:g6:rocks-and-minerals:1}
\textbf{1.} Graniet, een \omterm{def:g6:rocks-and-minerals:rock-families}{stollingsgesteente}; de \omterm{def:g6:rocks-and-minerals:mineral}{mineralen} zijn veldspaat
(roze of wit), kwarts (grijs, glazig) en mica (zwart, glimmend).

\textbf{2.} Kalksteen, een \omterm{def:g6:rocks-and-minerals:rock-families}{afzettingsgesteente}. De schelpen waren van
dieren die in de \omterm{def:g2:water-on-earth:ocean}{zee} leefden; hun resten werden neergelegd en verhardden
tot het \omterm{def:g3:stones-pebbles-sand:rock}{gesteente}.

\textbf{3.} Leisteen, een \omterm{def:g6:rocks-and-minerals:rock-families}{metamorf gesteente}, ontstaan uit schalie.

\textbf{4.} Marmer, een \omterm{def:g6:rocks-and-minerals:rock-families}{metamorf gesteente}, ontstaan uit kalksteen.

\textbf{5.} Kwarts: niet gekrast door staal (ongeveer 6.5), het krast glas
(ongeveer 5.5); de \omterm{def:g6:rocks-and-minerals:hardness}{hardheid} is 7.

\textbf{6.} Gekrast door staal (ongeveer 6.5) maar niet door glas
(ongeveer 5.5): tussen 5.5 en 6.5. Het is veldspaat (orthoklaas, 6).

\textbf{7.} Mica (\omterm{def:g6:rocks-and-minerals:hardness}{hardheid} 2 tot 3).

\textbf{8.} Ja: de bestrating is kalksteen, gemaakt van calciet (3), dat
zachter is dan een koperen munt (ongeveer 3.5).

\textbf{9.} \omterm{def:g5:water-shapes-the-land:weathering}{Verwering}: de \omterm{def:g4:clouds-rain-snow:precipitation}{regen} slijt het oppervlak van de steen langzaam
af.

\textbf{10.} De \omterm{def:g1:landscapes-around-us:landform}{berg} verweert; de kwartskorrels worden door \omterm{def:g1:landscapes-around-us:water}{rivieren}
weggevoerd en in lagen neergelegd, op de bodem van een \omterm{def:g1:landscapes-around-us:water}{meer} of van de \omterm{def:g2:water-on-earth:ocean}{zee};
begraven en verhard kunnen ze een zandsteen worden, een
\omterm{def:g6:rocks-and-minerals:rock-families}{afzettingsgesteente}.

\textbf{11.} Alle drie de families zijn er: \omterm{def:g6:rocks-and-minerals:rock-families}{stollingsgesteente} (het graniet
van de fontein), \omterm{def:g6:rocks-and-minerals:rock-families}{afzettingsgesteente} (de kalksteen van de bestrating),
\omterm{def:g6:rocks-and-minerals:rock-families}{metamorf gesteente} (de leisteen van het dak en het marmer van het beeld).

\textbf{12.} Het hardste \omterm{def:g6:rocks-and-minerals:mineral}{mineraal} is het kwarts van het graniet: \omterm{def:g6:rocks-and-minerals:hardness}{hardheid}
\textbf{7} op de schaal van Mohs.
\end{solution}
